\section{Transformation Methods}
\label{sec:2.1}

Let $X \sim f$ be a random variable taking values on a set $E \subset \R^d$ and let $U \sim \Unif(0,1)^d$ be a
random variable uniformly distributed on the $d$-dimensional unit cube $(0,1)^d$. In this section, we
seek to find a map $\mathcal{T} : (0,1)^d \to E$ so that the random variable $\mathcal{T}(U)$ has the same distribution
as $X$, written $X \overset{d}{=} \mathcal{T}(U)$. Once we find such a \textit{push-forward} or \textit{transport} map $\mathcal{T}$, we can obtain
a sample from the target distribution $f$ by:
\begin{enumerate}
  \item Sampling $U^{(1)}, \ldots, U^{(N)} \iid \Unif(0,1)^d$; and
  \item Setting $X^{(n)} := \mathcal{T}(U^{(n)})$ for $1 \le n \le N$.
\end{enumerate}

It is natural to ask if such a transport map $\mathcal{T}$ is guaranteed to exist. We will provide an affirmative
answer and an explicit construction in this section under the assumption that the target admits
a positive p.d.f. For exposition purposes, we consider first the scalar case $d = 1$ in Subsection
\ref{sec:2.1.1}, and then generalize to the $d$-dimensional case in Subsection~\ref{sec:2.1.2}.

\subsection{Inverse Transformation Method}
\label{sec:2.1.1}

Let $X \sim f$ be a real-valued random variable, and denote by $F$ its c.d.f. The inverse transformation
method relies on the (generalized) inverse c.d.f.\ of $X$, given by
\[
  F^-(u) := \inf\{x : F(x) \ge u\},
\]
to transform a uniform sample into a sample from $f$. The following result shows that the inverse
c.d.f.\ is a valid transport map.

\begin{theorem}[Transport: Uniform to Target]
\label{thm:2.1}
Let $X$ be a real-valued random variable with inverse c.d.f.\ $F^-$ and let $U \sim \Unif(0,1)$. It holds that $F^-(U) \overset{d}{=} X$.
\end{theorem}

\begin{proof}
We only prove here the case where $F$ is a bijection from $\R$ to $(0,1)$. In such a case, it
holds that $F^- = F^{-1}$, as illustrated in Figure~\ref{fig:2.1}. Then, for any $x \in \R$,
\[
  \Prob\bigl(F^-(U) \le x\bigr) = \Prob\bigl(F^{-1}(U) \le x\bigr) = \Prob\bigl(U \le F(x)\bigr) = F(x),
\]
which shows that $F$ is the c.d.f.\ of the random variable $F^-(U)$. The general case where the c.d.f.\
$F$ may not be invertible is the content of Exercise~2.2.
\end{proof}

\begin{figure}[htbp]
\centering
\includegraphics[width=0.75\textwidth]{figures/fig_02_1}
\caption{\textit{A strictly increasing c.d.f.\ and associated inverse c.d.f.}}
\label{fig:2.1}
\end{figure}

The inverse transformation method produces a sample from the target distribution as follows:

\begin{algorithm}[H]
\caption{Inverse Transformation Method}
\label{alg:2.1}
\begin{algorithmic}[1]
\Require Target distribution $f$ with inverse c.d.f.\ $F^-$, sample size $N$.
\State Set $X^{(n)} := F^-(U^{(n)})$, \quad $1 \le n \le N$, \quad where $U^{(1)}, \ldots, U^{(N)} \iid \Unif(0,1)$.
\Ensure Sample $\{X^{(n)}\}_{n=1}^N \iid f$.
\end{algorithmic}
\end{algorithm}

\medskip
\noindent\fbox{\parbox{\dimexpr\linewidth-2\fboxsep-2\fboxrule}{%
\textbf{Pros and Cons}: The inverse transformation method is very efficient whenever $F^-$ has a
closed form expression that can be evaluated cheaply. However, Algorithm~\ref{alg:2.1} is only applicable to sample real-valued random variables.}}
\medskip

\begin{example}[Inverse Transformation Method: Exponential Distribution]
\label{ex:2.2}
Suppose that $X \sim \Exp(\lambda)$, so that $F(x) = 1 - e^{-\lambda x}$, $x \ge 0$. Let $U \sim \Unif(0,1)$. Then,
\[
  -\frac{1}{\lambda}\log(1 - U) \sim \Exp(\lambda).
\]
Using that $1 - U \overset{d}{=} U$, we also have that
\[
  -\frac{1}{\lambda}\log(U) \sim \Exp(\lambda).
\]
Figure~\ref{fig:2.2} shows a histogram of $N = 10^5$ draws generated from an $\Exp(1)$ distribution
using the inverse transformation method.
\end{example}

\begin{figure}[htbp]
\centering
\includegraphics[width=0.75\textwidth]{figures/fig_02_2}
\caption{\textit{Inverse transformation method for sampling from an Exponential(1) distribution.
Uniform samples (first row) are transformed into samples from the target distribution (second row).}}
\label{fig:2.2}
\end{figure}

So far, we have discussed how to transform a uniform sample into a sample from a given
target distribution. In some applications, it is important to use a proposal distribution that is not
uniform. The following theorem provides a transport map from proposal to target.

\begin{theorem}[Transport: Proposal to Target]
\label{thm:2.3}
Let $Z$ be a real-valued random variable with continuous and strictly increasing c.d.f.\ $F_Z$ and let $X$ be a real-valued random variable with inverse c.d.f.\ $F_X^-$. Then, it holds that $F_X^- \circ F_Z(Z) \overset{d}{=} X$.
\end{theorem}

\begin{proof}
By Theorem~\ref{thm:2.1}, it suffices to show that $F_Z(Z) \sim \Unif(0,1)$. To that end, note that, for
any $z \in (0,1)$,
\[
  \Prob\bigl(F_Z(Z) \le z\bigr) = \Prob\bigl(Z \le F_Z^{-1}(z)\bigr) = F_Z\bigl(F_Z^{-1}(z)\bigr) = z,
\]
as desired.
\end{proof}

Let $X \sim f$ and let $Z \sim g$, where $g$ is a strictly positive p.d.f. In view of Theorem~\ref{thm:2.3}, we
can obtain a draw from $f$ by (1) drawing $Z^{(n)} \sim g$; and (2) setting $X^{(n)} := F_X^- \circ F_Z(Z^{(n)})$.

\subsection{Knothe-Rosenblatt Rearrangement}
\label{sec:2.1.2}

We now generalize the inverse transformation method to target distributions on $\R^d$, $d \in \N$.
For ease of exposition, let us consider first the case $d = 2$. Let $X = (X_1, X_2) \sim f$ and let
$U = (U_1, U_2) \sim \Unif(0,1)^2$. The idea is to \textit{disintegrate} the target $f$ as the product of the
marginal of $X_1$ times the conditional of $X_2$ given $X_1$:
\[
  f(x) = f_{X_1}(x_1) f_{X_2|X_1=x_1}(x_2|x_1), \quad x = (x_1, x_2) \in \R^2.
\]

We then sequentially apply the inverse transformation method to the one-dimensional targets
$f_{X_1}$ and $f_{X_2|X_1=x_1}$ as will be detailed next.

Denote by $F_{X_1}$ the marginal c.d.f.\ of $X_1$. By the inverse transformation method, we know
that $F_{X_1}^-(U_1) \overset{d}{=} X_1$. Now, for any $x_1 \in \R$ denote by $F_{X_2|X_1=x_1}$ the c.d.f.\ of $X_2 | X_1 =
x_1$. By the inverse transformation method, we know that $F_{X_2|X_1=x_1}^-(U_2) \overset{d}{=} X_2|X_1 = x_1$. To
summarize, define recursively maps $\mathcal{T}^1$ and $\mathcal{T}^2$ by
\begin{align*}
  \mathcal{T}^1 : (0,1) \to \R, & \qquad \mathcal{T}^2 : (0,1)^2 \to \R, \\
  u_1 \mapsto F_{X_1}^-(u_1), & \qquad (u_1, u_2) \mapsto F_{X_2|X_1=\mathcal{T}^1(u_1)}^-(u_2).
\end{align*}

Then, the triangular map $\mathcal{T} : (0,1)^2 \to \R^2$ given by
\begin{equation}
  \mathcal{T}(u) = \begin{bmatrix} \mathcal{T}^1(u_1) \\ \mathcal{T}^2(u_1, u_2) \end{bmatrix}, \qquad u = (u_1, u_2) \in \R^2,
  \label{eq:2.1}
\end{equation}
satisfies that $\mathcal{T}(U) \overset{d}{=} X$. The map $\mathcal{T}$ defined in \eqref{eq:2.1} is called the \textit{Knothe-Rosenblatt rearrangement} from $\Unif(0,1)^2$ to $f$. The construction can be generalized to higher dimensions:

\begin{theorem}[Transport via Triangular Maps]
\label{thm:2.4}
Let $f$ be a strictly positive p.d.f.\ on $\R^d$. Let $X \sim f$ and let $U \sim \Unif(0,1)^d$. For $1 \le i \le d$, define maps $\mathcal{T}^i$ as follows:\footnote{For vector $u = (u_1, \ldots, u_d) \in \R^d$ and $1 \le i \le d$, we denote by $u_{1:i} := (u_1, \ldots, u_i) \in \R^i$ the vector defined by the first $i$ coordinates of $u$. Likewise, for triangular map $\mathcal{T} : \R^d \to \R^d$, we denote by $\mathcal{T}^{1:i} : \R^i \to \R^i$ the map defined by the first $i$ components of $\mathcal{T}$.}
\begin{align}
  \mathcal{T}^1(u_1) &= F_{X_1}^-(u_1), \label{eq:2.2} \\
  \mathcal{T}^i(u_{1:i}) &:= F_{X_i|X_{1:i-1}=\mathcal{T}^{1:i-1}(u_{1:i-1})}^-(u_i), \qquad 2 \le i \le d. \label{eq:2.3}
\end{align}
Then, the triangular map
\begin{equation}
  \mathcal{T}(u) = \begin{bmatrix} \mathcal{T}^1(u_1) \\ \mathcal{T}^2(u_1, u_2) \\ \vdots \\ \mathcal{T}^d(u_1, \ldots, u_d) \end{bmatrix}, \qquad u = (u_1, \ldots, u_d) \in \R^d,
  \label{eq:2.4}
\end{equation}
satisfies that $X \overset{d}{=} \mathcal{T}(U)$. Moreover, for all $1 \le i \le d$, the components $\mathcal{T}^i : (0,1)^d \to \R$ are monotonically increasing in the $i$-th variable.
\end{theorem}

The map $\mathcal{T}$ defined in Theorem~\ref{thm:2.4} is called the Knothe-Rosenblatt (KR) rearrangement
from $\Unif(0,1)^d$ to $f$. The map $\mathcal{T}$ is referred to as a \textit{triangular} map since its Jacobian is a
(lower) triangular matrix. The KR rearrangement can be used for sampling, as summarized in
the following algorithm:

\begin{algorithm}[H]
\caption{Sampling via the KR Rearrangement}
\label{alg:2.2}
\begin{algorithmic}[1]
\Require KR rearrangement $\mathcal{T}$ from $\Unif(0,1)^d$ to $f$, sample size $N$.
\State Set $X^{(n)} := \mathcal{T}(U^{(n)})$, \quad $1 \le n \le N$, \quad where $U_1, \ldots, U^{(N)} \iid \Unif(0,1)^d$.
\Ensure Sample $\{X^{(n)}\}_{n=1}^N \iid f$.
\end{algorithmic}
\end{algorithm}

\medskip
\noindent\fbox{\parbox{\dimexpr\linewidth-2\fboxsep-2\fboxrule}{%
\textbf{Pros and Cons}: As with other transformation methods, sampling via the KR rearrangement is
very efficient if we can find and evaluate the KR map. However, computing the KR rearrangement can be expensive in high dimensions. Conditional independence of the target-reference pair leads to sparsity in the KR map that can be exploited when computing or estimating the KR map. We refer to Section~\ref{sec:2.3} for further discussion on this topic.}}
\medskip

Using Theorem~\ref{thm:2.3}, it is possible to generalize the construction of the KR rearrangement to
transport a strictly positive proposal p.d.f.\ $g$ on $\R^d$ into a strictly positive target p.d.f.\ $f$. This is
the content of Exercise~2.4.
